\section{The New Intelligence Allocation Problem}

This vision responds to two changes reshaping the problem that AI systems must
solve. The demand for intelligence is becoming dynamic as applications move from isolated
requests to agentic trajectories. At the same time, the supply of intelligence
is becoming heterogeneous as models diverge in capability, specialization, and
operating cost. The opportunity---and the systems challenge---is to match the
two continuously.

\subsection{From requests to trajectories}

The familiar model interface was designed around a bounded exchange: a request
arrives, a model produces a response, and the interaction ends. Within that
frame, it is reasonable to characterize the request once, select a model, and
evaluate the resulting answer.

Agents change the unit of work. A long-running task develops through a sequence
of observations, hypotheses, actions, artifacts, failures, and verifications.
Each transition changes what the system knows and what it can responsibly do
next. The intelligence required at the beginning of a trajectory may bear
little resemblance to the intelligence required near its end.

The difference is not merely that some steps are harder than others. The
\emph{kind} of capability required also changes. One step may need broad
exploration, another domain knowledge, another precise construction, another
an adversarial critique, and another a deterministic observation from the
environment. A consequential action may require stronger judgment even when
the underlying reasoning is simple.

Most importantly, this demand cannot be fully known in advance. A tool result
can eliminate the apparent hard part of a task. A failed test can reveal a
deeper problem that was invisible in the original request. A candidate answer
can turn a vague concern into a falsifiable claim. Intelligence demand in an
agentic workflow is therefore state-dependent, multidimensional, and
\emph{partially observable}: it is progressively revealed by doing the work.

\subsection{From a model hierarchy to a capability portfolio}

While demand is becoming dynamic, supply is becoming more diverse. Models
differ because of training data, parameter scale, architecture, post-training,
tool exposure, domain specialization, and deployment surface. These differences
produce distinct strengths and failure modes across reasoning, code, science,
search, critique, tool use, latency, reliability, privacy, and cost.

The result is not simply a longer leaderboard. Model capability is conditional.
The best standalone answerer may not be the best critic. A smaller model may be
more valuable when it searches a different region, supplies specialist
knowledge, or performs a bounded check. Two frontier models may be individually
strong yet jointly redundant because their errors are correlated.

Cost is part of this capability landscape, not an unrelated commercial filter.
It determines how broadly a system can search, how frequently it can verify,
and how many real tasks it can learn from. A model that is useful only under an
unbounded inference budget has a different system capability from one whose
strength can be deployed repeatedly and selectively.

AI systems are therefore crossing two boundaries at once:

\principle{The unit of work is becoming a trajectory. The supply of
intelligence is becoming a portfolio.}

\subsection{Why routing is not enough}

A model router is the natural bridge between heterogeneous demand and
heterogeneous supply. For a bounded request whose requirements are visible at
the outset, routing can be highly effective: classify the request, choose an
appropriate model, and escalate when a cheaper model is unlikely to suffice.

The limitation appears when routing is extended from bounded requests to an
agentic trajectory. A router can run again at every turn, but each decision
still classifies the current demand before acquiring the evidence that may
reveal it. It assumes that the requirement can be reduced to a category or
difficulty level and that one model can own the resulting step. In a changing
workflow, these assumptions do not hold reliably.

Determining difficulty can itself require much of the reasoning needed to solve
the problem. More fundamentally, difficulty is the wrong control variable. A
system can know that a step is difficult without knowing whether it needs
another solver, a domain specialist, an external observation, an independent
critic, stronger synthesis, or no further computation at all. What matters is
not only how much intelligence the task requires, but what remains unresolved
and what kind of evidence could resolve it.

This is the \emph{intelligence allocation problem}: how to match a
partially-observed, continuously changing demand for intelligence with a
heterogeneous and economically non-uniform supply of capability.

Solving it requires a different abstraction. Instead of assigning each step to
one model before its uncertainty has been explored, the system must repeatedly
decide what computation is valuable next. Instead of treating model outputs as
final answers, it must use them to update its understanding of the task. And
instead of assuming that more compute is always better, it must know when the
evidence is sufficient to act.
